% Generated draft: requires mathematical and source review.
\documentclass[11pt]{article}
\usepackage[margin=25mm]{geometry}
\usepackage{fontspec}
\setmainfont{Latin Modern Roman}
\usepackage{amsmath,amssymb,mathtools,unicode-math}
\setmathfont{Latin Modern Math}
\usepackage{enumitem,needspace,xurl,hyperref}
\setlength{\parindent}{0pt}
\setlength{\parskip}{5pt}
\setlength{\emergencystretch}{3em}
\newcommand{\N}{\mathbb{N}}
\newcommand{\Z}{\mathbb{Z}}
\newcommand{\Q}{\mathbb{Q}}
\newcommand{\R}{\mathbb{R}}
\newcommand{\C}{\mathbb{C}}
\newcommand{\F}{\mathbb{F}}
\newcommand{\A}{\mathbb{A}}
\newcommand{\PP}{\mathbb{P}}
\newcommand{\E}{\mathbb{E}}
\newcommand{\eps}{\varepsilon}
\newcommand{\dd}{\,\mathrm{d}}
\begin{document}
\section*{UnsolvedMath ID: 30003169}
\textbf{Source problem code:} catalogue entry
\subsection{Definitions and mathematical statement}
Let \(f\), \(K\), and \(c\) denote the underlying function, kernel, and level appearing in the supplied record, and let \(\widehat f_h\) denote its kernel estimator with bandwidth \(h>0\). In the ambient space \(\R^d\), define
\[
L=\{x\in\R^d:f(x)\geq c\},\qquad
\widehat L_h=\{x\in\R^d:\widehat f_h(x)\geq c\}.
\]
The source writes the mean distance-in-measure risk as
\[
\operatorname{MDM}(h)=I_1(h)+I_2(h).
\]
Here the dependence of the estimator and risk terms on the sample size \(n\) is suppressed in the source notation.

The proposed question is whether there exists a constant \(B=B(f,K,c)\) such that, writing \(h=h_n\) for the bandwidth at sample size \(n\),
\[
\lim_{n\to\infty}h_n^{-2}I_2(h_n)=B(f,K,c),
\]
without assuming
\[
nh_n^{d+4}\longrightarrow0.
\]
\textbf{Unresolved qualifications.} The supplied record does not define \(I_1\) or \(I_2\), specify the measure and expectation defining \(\operatorname{MDM}\), or give the sampling model and kernel-estimator formula. It also omits the admissible classes of \(f\), \(K\), and \(c\), boundary regularity assumptions, and the remaining conditions and quantifiers on bandwidth sequences. Consequently, the displayed decomposition does not determine \(I_2\), and this is a transcription of the intended limit question rather than a fully specified mathematical conjecture. Removing the displayed bandwidth condition leaves the other intended hypotheses unresolved.
\subsection{Short English statement}
The question asks whether the term called \(I_2\) in a kernel level-set estimation risk has an \(h^2\)-scaled limit depending only on \(f\), \(K\), and \(c\), without the condition \(nh^{d+4}\to0\). The risk decomposition and remaining hypotheses are missing, so a precise conjecture requires further source clarification.
\subsection{Sources}
\begin{itemize}
\item Local registry citation: supplied complete UnsolvedMath source record, permanent ID \(30003169\), title ``Second-Order Bias in Kernel Level-Set Estimation,'' Oberwolfach Reports Open Problems collection, exact source problem code \texttt{catalogue entry}. No named proposer is supplied; the proposed year is recorded as 2016.
\item Supplied report citation: \emph{Statistics for Shape and Geometric Features (2017)}, Oberwolfach Reports, pp.\ 1821--1874, DOI \texttt{10.4171/owr/2016/32}. \url{https://doi.org/10.4171/owr/2016/32}. The background attributes the collection to workshop open-problem discussions documented in Oberwolfach Reports. The same URL appears in the literature sources under ``Oberwolfach Report on density-level-set risk.''
\item Supplied literature reference: Qiao, ``Asymptotics and optimal bandwidth.'' \url{https://arxiv.org/abs/1707.09697}.
\item Supplied literature assessment, checked on \texttt{2026-08-21}: status ``partially solved.'' The assessment states that Qiao proves asymptotic risk expansions and derives an optimal bandwidth under explicit regularity and bandwidth assumptions, but that the unrestricted limit without the extra bandwidth condition is not established by that result. This is supplied dated curation derived from the OWR record, not an independently verified literature assessment. The background explicitly requests independent verification of time-sensitive status and cited claims.
\end{itemize}
\end{document}
